\documentclass[11pt,a4paper]{article}

\usepackage[margin=1in]{geometry}
\usepackage{amsmath,amssymb,amsthm,bm}
\usepackage[dvipsnames]{xcolor}
\usepackage{enumitem}
\usepackage{booktabs}
\usepackage{array}
\usepackage{hyperref}
\usepackage{url}
\makeatletter
\g@addto@macro{\UrlBreaks}{\do\.\do\_\do\-}
\makeatother
\newcommand{\lmod}[2]{\href{http://www.lemma.cn/lean/?module=#2}{\nolinkurl{#1}}}

\hypersetup{
  colorlinks=true,
  linkcolor=blue,
  citecolor=blue,
  urlcolor=blue
}

\newtheorem{theorem}{Theorem}[section]
\newtheorem{lemma}[theorem]{Lemma}
\newtheorem{proposition}[theorem]{Proposition}
\newtheorem{definition}[theorem]{Definition}
\newtheorem{remark}[theorem]{Remark}

\newcommand{\R}{\mathbb{R}}
\newcommand{\N}{\mathbb{N}}
\newcommand{\E}{\mathbb{E}}
\newcommand{\Prb}{\mathbb{P}}
% colour convention of the lemma.cn renderer
\newcommand{\rv}[1]{{\color{red}#1}}       % random variable
\newcommand{\ra}[1]{{\color{magenta}#1}}   % random argument
\newcommand{\sv}{\rv{s}}
\newcommand{\av}{\rv{a}}
\newcommand{\rw}{\rv{r}}
\newcommand{\grad}{\nabla_{\theta}}
\newcommand{\tsum}{\textstyle\sum'}
\newcommand{\lse}{\operatorname{logsumexp}}
\newcommand{\Y}{\mathcal{Y}}
\emergencystretch=3em

\title{A Lean~4 Perspective on NLP:\\ from CRF to PPO}
\author{LiZhi Zhou\\[0.3em]
\small Project artifact: \href{https://github.com/math-proof/lemma}{math-proof/lemma}\\[0.4em]
\small
Markov product:
\href{http://www.lemma.cn/lean/?module=Random.Prob.eq.Mul_Prod_MulProbS.of.Ne_0.Eq.Eq.Eq}{$P_t$},
\quad
forward recursion, NLL:
\href{http://www.lemma.cn/lean/?module=Tensor.Imp_Eq_Add_LogSumExp.EqNegLogProb.of.Ne0ProbJoint.Indep.y_given_x}{$\log Z-\text{score}$},
\quad
Viterbi:
\href{http://www.lemma.cn/lean/?module=Tensor.And.of.Ne_0.Eq.Eq.Eq}{$\max$}\\[0.25em]
forward--backward:
\href{http://www.lemma.cn/lean/?module=Tensor.And.of.All_Eq_Sum_Mul.All_Eq_1.All_Eq_Mul.All_Eq.forward_backward}{$z_t\beta_t$},
\quad
Bellman:
\href{http://www.lemma.cn/lean/?module=Tensor.And.Eq.Expect.of.Eq_Conditioned.Eq_Expect.Eq_Expect.Bellman}{$V,Q$},
\quad
REINFORCE:
\href{http://www.lemma.cn/lean/?module=Tensor.Eq.Dot.Grad.Expect.of.Eq_Conditioned.IsFinite.policy_gradient_theorem}{Lean~4}\\[0.25em]
Unbiased advantage estimate:
\href{http://www.lemma.cn/lean/?module=Tensor.EqDot_GradExpect.of.Eq_Conditioned.Eq_Expect.IsFinite.IsFinite.unbiased_advantage_estimate}{Lean~4},
\quad
generalized advantage estimate:
\href{http://www.lemma.cn/lean/?module=Tensor.EqDot_GradExpect.of.Eq_Conditioned.Eq_Expect.IsFinite.IsFinite.generalized_advantage_estimate}{Lean~4}}
\date{}

\begin{document}
\maketitle

\begin{abstract}
Linear-chain conditional random fields (CRFs) for sequence labeling and proximal policy optimization (PPO) for reinforcement learning from feedback look like unrelated corners of NLP.
We organize a Lean~4 account of the path between them around three results: the forward and backward propagation of the linear-chain CRF, the Bellman equation, and the policy gradient theorem with its REINFORCE and generalized-advantage-estimation (GAE) forms.
All three rest on one device: a Markov structure turns a sum over exponentially many paths (or over an infinite horizon) into a one-step recursion, run forward for the CRF partition function and backward for suffix sums and for the value function.
We say exactly what is machine-checked in Lean~4 with mathlib.
For the CRF we formalize the Markov product, the decomposition of the log-score into transition and emission terms, the forward recursion in log-sum-exp form with $-\log p(y\mid x)=\log Z(x)-\text{score}(x,y)$, the Viterbi max-plus recursion (maximal value only), and a forward--backward theorem for an arbitrary chain of real weights: the forward recursion, the backward recursion, the identity $\sum_{ys:\,ys_t=a}P_m(ys)=z_t(a)\,\beta_t(a)$ for the total weight of the paths through label $a$ at time $t$, and the independence of $Z$ from the cut $t$.
For reinforcement learning we formalize a discounted Markov decision process (MDP) with finite state and action spaces and its trajectory measure, the Bellman equations, the policy gradient theorem by induction on the horizon in truncated, action-value and REINFORCE form, the zero-mean score identity that makes a baseline vanish, the unbiased discounted advantage, and GAE with the exact value function for every $\lambda\in[0,1]$.
Not formalized, and labeled as such, are the normalized posterior and pairwise marginals, the identity $\nabla\log Z=\E[\nabla\,\text{score}]$, the Viterbi argmax, and every ingredient of the PPO algorithm (clipped surrogate, importance ratio, KL and value clipping, learned critic); the analogy between the forward, backward and Bellman recursions, and the comparison of sequence labeling with sequence-to-sequence generation, are conceptual and are not Lean theorems.
All checked theorems depend only on the axioms \texttt{propext}, \texttt{Classical.choice} and \texttt{Quot.sound}.
\end{abstract}

\section{Introduction}
\label{sec:intro}

Sequence labeling with a linear-chain CRF~\cite{lafferty2001crf,sutton2012crf} and reinforcement learning from human feedback with PPO~\cite{schulman2017ppo,ziegler2019finetuning,ouyang2022instructgpt} are usually taught in different chapters.
The first is supervised: an observed sentence $x$ is mapped to a tag sequence $y$ of the same length by maximizing $\log p(y\mid x)$.
The second is a reinforcement-learning algorithm: a language model generates tokens, a scalar reward judges the finished text, and the policy is updated by a clipped policy-gradient step.
This paper is about what they share, seen through what we have been able to prove in Lean~4~\cite{moura2021lean4} on top of mathlib~\cite{mathlib2020,mathlib}, in the \texttt{lemma} library~\cite{lemma_repo}.

\paragraph{Three core topics.}
The paper is organized around three results, each with its own section and its own Lean statements.
\begin{enumerate}[leftmargin=1.6em,itemsep=0.25em]
  \item \textbf{Forward and backward propagation of the linear-chain CRF} (Section~\ref{sec:crf}).
  The forward recursion $z_{t+1}(a)=\sum_bz_t(b)\,\psi_{t+1}(b,a)$ computes the partition function $Z$ in $O(mk^2)$ operations instead of summing $k^{m+1}$ label paths; the backward recursion $\beta_t(a)=\sum_c\psi_{t+1}(a,c)\,\beta_{t+1}(c)$ summarizes the suffix, and $z_t(a)\beta_t(a)$ is the total weight of the label paths through $a$ at time $t$, the numerator of the posterior marginal.
  Lean has the forward recursion with $-\log p(y\mid x)=\log Z-\text{score}$ for the hidden-Markov-conditional case, the Viterbi value recursion, and a forward--backward theorem (Theorem~\ref{thm:fb}) for an arbitrary chain of weights.
  \item \textbf{The Bellman equation} (Section~\ref{sec:bellman}).
  For a discounted MDP with finite state and action spaces, the value function satisfies $V_t(s_t)=\E[\rw_t+\gamma V_{t+1}(\ra{s}_{t+1})\mid\sv_t=s_t]$, and likewise for the action value $Q_t$; Lean proves this for the trajectory measure of the MDP.
  \item \textbf{The policy gradient theorem, REINFORCE and GAE} (Section~\ref{sec:pg}).
  $\grad\E[\tsum_t\gamma^t\rw_t]=\E[\tsum_t\gamma^tG_t\grad\log\pi_\theta(\av_t\mid\sv_t)]$ is proved by induction on the horizon, with a finite-horizon identity that has an explicit remainder; replacing $G_t$ by the unbiased advantage, or by the generalized advantage $\hat A^\lambda_t$ built from the exact value function, leaves the gradient unchanged.
\end{enumerate}
Section~\ref{sec:ppo} says which parts of PPO lie outside these statements, and Section~\ref{sec:gen} is a short conceptual comparison of sequence labeling with sequence-to-sequence generation that is not a theorem.

\paragraph{The common thread.}
In every model of this paper the probability of a length-$n$ object is a product of one-step factors (the Markov assumption), so its logarithm is a sum of per-step terms: the CRF score, the log-probability of a trajectory, the log-likelihood of a sentence under a language model.
Two consequences are used.
First, sums over exponentially many paths are computed by a one-step recursion: the forward and backward recursions of the CRF are a sum-product dynamic program (and Viterbi its max-sum version), and the Bellman recursion is an expectation recursion of the same shape, run backward in time (Section~\ref{sec:analogy}).
We present this as an analogy, not as a theorem: no Lean statement relates the CRF recursions to the Bellman recursion.
Second, when each step factor sums to one, $\sum_a p(a)=1$ implies $\E_{a\sim p}[\grad\log p(a)]=0$, which is why a baseline can be subtracted from a policy gradient without changing it (Section~\ref{sec:zeroscore}).

\paragraph{Contributions.}
\begin{enumerate}[leftmargin=1.6em,itemsep=0.25em]
  \item CRF propagation (Sections~\ref{sec:markov} and~\ref{sec:crf}): the Markov product \lmod{crf.markov}{Random.Prob.eq.Mul_Prod_MulProbS.of.Ne_0.Eq.Eq.Eq}, the log-score \lmod{crf.markov.logits}{Tensor.Eq.of.Ne_0.Eq.Eq.Eq.Eq_log.Eq_log.Eq_log} and \lmod{crf.logits}{Tensor.Imp.of.Eq}, the forward recursion with $-\log p(y\mid x)=\log Z-\text{score}$, \lmod{y_given_x}{Tensor.Imp_Eq_Add_LogSumExp.EqNegLogProb.of.Ne0ProbJoint.Indep.y_given_x}, the Viterbi value recursion \lmod{crf.viterbi}{Tensor.And.of.Ne_0.Eq.Eq.Eq}, and the new forward--backward theorem \lmod{crf.forward_backward}{Tensor.And.of.All_Eq_Sum_Mul.All_Eq_1.All_Eq_Mul.All_Eq.forward_backward}, whose hypotheses are the one-step prefix factorization and the backward recursion (Theorem~\ref{thm:fb}).
  \item The Bellman equation \lmod{Bellman}{Tensor.And.Eq.Expect.of.Eq_Conditioned.Eq_Expect.Eq_Expect.Bellman} on a trajectory-level MDP whose trajectory law has a proved Markov property (\texttt{joint\_succ}, \texttt{hist\_step}) (Section~\ref{sec:bellman}).
  \item The policy gradient theorem by induction on the horizon, in truncated, action-value and REINFORCE form (\lmod{policy_gradient_theorem}{Tensor.Eq.Dot.Grad.Expect.of.Eq_Conditioned.IsFinite.policy_gradient_theorem}), the unbiased advantage estimate \lmod{unbiased_advantage_estimate}{Tensor.EqDot_GradExpect.of.Eq_Conditioned.Eq_Expect.IsFinite.IsFinite.unbiased_advantage_estimate} and GAE with the exact value function for all $\lambda\in[0,1]$, \lmod{generalized_advantage_estimate}{Tensor.EqDot_GradExpect.of.Eq_Conditioned.Eq_Expect.IsFinite.IsFinite.generalized_advantage_estimate}, with its weighted-average form (Section~\ref{sec:pg}).
  \item An explicit statement of what is \emph{not} covered, in particular the ingredients of PPO (Sections~\ref{sec:ppo} and~\ref{sec:discussion}), and two comparison tables (Section~\ref{sec:tables}).
\end{enumerate}
The textbook syntactic sugar of the library for limits, expectations and probabilities, which makes the Lean statements read like the formulas of the paper, is summarised in Section~\ref{sec:notation}.

\medskip
\noindent\fbox{\parbox{\dimexpr\linewidth-2\fboxsep-2\fboxrule\relax}{\small
\textbf{Honest scope.}
(i) The downstream log-domain CRF lemmas (logits, forward recursion with the negative log-likelihood, Viterbi) take the one-step prefix factorization as a \emph{hypothesis}; for the probabilistic form it is derived from the emission and first-order Markov conditional independences (Theorem~\ref{thm:markov}), but it is not derived from a concrete construction of the model.
(ii) The backward recursion is formalized only inside the forward--backward theorem (Theorem~\ref{thm:fb}), which is a statement about arbitrary real weights: it does not define the CRF distribution $p(y\mid x)=\prod_t\psi_t/Z$ as a model, and the normalized posterior marginals, the pairwise marginals, the identity $\nabla\log Z=\E[\nabla\,\text{score}]$ and the Viterbi argmax with back-pointers are \emph{not} formalized.
(iii) There is no sequence-to-sequence, teacher-forcing, exposure-bias or label-bias lemma.
(iv) The MDP is formalized; PPO is not: there is no clipped surrogate, importance ratio, KL penalty, value clipping, learned critic, or monotone-improvement or convergence statement.
(v) GAE is proved only with the exact value function, for finite state and action spaces, bounded rewards and $\gamma<1$.
Sections~\ref{sec:analogy}, \ref{sec:gen} and~\ref{sec:tables} contain conceptual comparisons and are labeled as such.}}

\section{Background and related work}
\label{sec:related}

\paragraph{HMMs, MEMMs and CRFs.}
A hidden Markov model (HMM)~\cite{rabiner1989hmm} is generative: it models the joint $p(x,y)=\prod_tp(y_t\mid y_{t-1})\,p(x_t\mid y_t)$ with locally normalized factors.
Maximum-entropy Markov models (MEMMs)~\cite{mccallum2000memm} are discriminative but still locally normalized, $\prod_tp(y_t\mid y_{t-1},x)$, and suffer from the label bias problem; conditional random fields~\cite{lafferty2001crf} remove it by normalizing globally with one partition function $Z(x)$ over whole label sequences.
Sutton and McCallum~\cite{sutton2012crf} give the standard tutorial treatment, including the forward, backward and Viterbi recursions.
Neural CRFs put a neural network under the unary scores: Collobert et al.~\cite{collobert2011nlp} train a sentence-level likelihood with transition scores, Huang et al.~\cite{huang2015bilstmcrf}, Lample et al.~\cite{lample2016ner} and Ma and Hovy~\cite{ma2016bilstmcnnscrf} use a BiLSTM encoder with a CRF layer, and Andor et al.~\cite{andor2016global} argue for globally normalized transition-based networks.
Su Jianlin's blog gives a concise derivation in Chinese of the CRF loss and its recursive normalizer~\cite{su2017crf}; further CRF-related posts are listed on the site's search page~\cite{su2018search}. GlobalPointer is a span-based alternative to CRF for NER~\cite{su2022globalpointer}.

\paragraph{Sequence-to-sequence models and exposure bias.}
Neural sequence-to-sequence models~\cite{sutskever2014seq2seq,bahdanau2015nmt,vaswani2017attention} factor $p(y\mid x)=\prod_tp(y_t\mid x,y_{<t})$ and are trained with the true prefix (teacher forcing) but decoded with their own; this mismatch is called exposure bias and is attacked by scheduled sampling~\cite{bengio2015scheduled}, sequence-level training~\cite{ranzato2016mixer} and beam-search optimization~\cite{wiseman2016bso}.
A small example in which beam search with a larger beam returns a better sequence is in Section~\ref{sec:gen}.

\paragraph{Policy gradients, GAE and PPO.}
The policy gradient theorem is due to Sutton et al.~\cite{sutton2000policy}, the likelihood-ratio estimator and the reinforcement baseline to Williams~\cite{williams1992simple}; Greensmith, Bartlett and Baxter~\cite{greensmith2004variance} analyse baselines as variance reducers.
Schulman et al.~\cite{schulman2016gae} introduced generalized advantage estimation, used there with trust region policy optimization (TRPO)~\cite{schulman2015trpo}; TRPO predates GAE and does not define it.
PPO~\cite{schulman2017ppo} replaces the trust-region constraint by a clipped surrogate objective and computes its advantages with a truncated (finite-horizon) GAE, its Eqs.~(10)--(11).
RLHF for language models~\cite{ziegler2019finetuning,ouyang2022instructgpt} runs PPO on tokens with a reward-model score and a KL penalty to a reference model, and GRPO~\cite{shao2024deepseekmath} is a PPO variant without a learned critic.
Standard references for finite MDPs are Puterman~\cite{puterman1994markov} and Sutton and Barto~\cite{sutton2018reinforcement}.

\paragraph{Formalizations.}
CertRL~\cite{vajjha2021certrl} formalizes convergence of value and policy iteration in Coq, and Chevallier and Fleuriot~\cite{chevallier2021isabelle} formalize finite MDPs with rewards in Isabelle/HOL, including the Bellman equation, the existence of an optimal policy for $\gamma<1$, and value and policy iteration; neither treats policy gradients.
In Lean~4, TorchLean~\cite{torchlean} proves structural and dynamic-programming facts about MDPs and Bellman operators and contains policy-gradient and PPO training code, but no policy gradient theorem.
Closest to Section~\ref{sec:pg} is the \texttt{PolicyGradient} module of Zhang's \texttt{rl-theory-in-lean}~\cite{zhang2026pg}, an experimental, largely machine-generated development (its commits are co-authored with an AI assistant) that was committed in March 2026 and later withdrawn from the project's main branch; it remains accessible at the cited commit.
For finite state and action spaces and a policy with finitely many real parameters it proves the occupancy-measure form
$\grad J(\theta)=(1-\gamma)^{-1}\sum_{s,a}d^{\pi_\theta}(s)\,\pi_\theta(a\mid s)\,Q^{\pi_\theta}(s,a)\,\grad\log\pi_\theta(a\mid s)$
via a Banach fixed point and a Neumann series, requires a strictly positive policy, and has no trajectories, returns or advantages.
Ours works at the level of trajectories: the law of a trajectory is mathlib's \texttt{Kernel.trajMeasure}~\cite{mathlib}, an implementation of the Ionescu-Tulcea theorem~\cite{kallenberg2002foundations}, the value functions are conditional expectations, and we prove the exact finite-horizon identity of the textbook derivation~\cite[\S13.2]{sutton2018reinforcement} by induction before passing to the limit.
To the best of our knowledge, based on searches of GitHub, arXiv and the web, this is the first Lean~4 proof that follows the textbook derivation; the two proofs share only the textbook starting point (the one-step recursion of~\cite{sutton2000policy} and the log-derivative trick), and ours was developed independently.
We are not aware of earlier Lean~4 formalizations of CRF forward or backward recursions.

\section{Lean notation in brief}
\label{sec:notation}

The statements below are written in the library's \emph{textbook syntactic sugar}: surface syntax that compresses verbose measure-theoretic or filter-theoretic expressions into near-textbook form while elaborating to complete, machine-checked definitions, with no new axioms (the generic term \emph{syntactic sugar} goes back to Landin~\cite{landin1964mechanical}).
We recall only what the later sections use.

\begin{center}
\small
\begin{tabular}{@{}>{\raggedright\arraybackslash}p{0.47\linewidth}>{\raggedright\arraybackslash}p{0.48\linewidth}@{}}
\toprule
Lean surface form & meaning \\
\midrule
\texttt{lim [n $\to$ $\infty$] e = a} & \texttt{Filter.Tendsto (fun n => e) atTop (nhds a)}; always the equation form, which asserts convergence \\
\texttt{$\mathbb{E}$[x: $\pi$](f x | y = y0)} & $\E_\pi[f(\rv{x})\mid\rv{y}=y_0]$, the expectation under mathlib's conditional measure; $0$ at a null event \\
\texttt{$\mathbb{P}$[$\pi$](x = x0 | y = y0)} & $\Prb_\pi(\rv{x}=x_0\mid\rv{y}=y_0)$ (densities w.r.t.\ counting measure give probability masses) \\
\texttt{x[:n]} & the random vector $(\rv{x}_0,\dots,\rv{x}_{n-1})$; for a plain sequence $f$, \texttt{f[:n]} is $(f_0,\dots,f_{n-1})$ \\
\texttt{($\mathbb{P}$[$\pi$](\dots) : $\mathbb{R}$).log} & $\log\Prb_\pi(\dots)$ (the scoped coercion $[0,\infty]\to\R$ is not a global instance) \\
\texttt{$\nabla$[$\theta$] e}, \texttt{sup[x] e $<$ $\infty$} & $\nabla_\theta e$ (Riesz representative of the Fr\'echet derivative); $e$ bounded above \\
\texttt{max[ys : Fin t $\to$ Y] f ys}, \texttt{v.exp.sum.log} & $\max_{ys}f(ys)$ over a finite index type; $\log\sum_b e^{v_b}$ (\texttt{v.exp}, \texttt{v.sum}, \texttt{v.max} act pointwise or reduce a vector) \\
\bottomrule
\end{tabular}
\end{center}

\paragraph{Colours.}
The renderer of lemma.cn typesets every statement in \LaTeX{} and colours identifiers by probabilistic role: \rv{red} for a random variable (a function on the sample space, the integrated variable of an $\E$ binder, the left side of $=$ in a $\Prb$ event), \ra{magenta} for a random argument (a random variable used as the argument of a function, so that the expression is itself random), and black for an observed value.
We use the same convention: $\sv_t,\av_t,\rw_t$ are red, $s_t,a_t,x,u,y$ are values, and $V_t(\ra{s}_t)$ is a function of the state evaluated along the trajectory.

\paragraph{Conditioning at null events and limits.}
Because $\mu[\,\cdot\mid B]=\mu(B)^{-1}\mu|_B$ is the zero measure when $\mu(B)=0$, a conditional expectation at an unreachable state is $0$; hypotheses on value functions are therefore restricted to \emph{reachable} values $s_t$, those with $\Prb_\theta(\sv_t=s_t)\ne0$.
Infinite sums are Lean's \texttt{tsum} ($\tsum$, the sum of a summable family and $0$ otherwise), and every limit is an explicit \texttt{Filter.Tendsto} statement.
The only limit lemma we state is the following, used for the remainder of the truncated policy gradient (Section~\ref{sec:pg}).

\begin{lemma}[\lmod{Lim.of.LtAbs.IsFinite}{Real.Eq_0.Lim.of.LtAbs.IsFinite}~\cite{lemma_lim}]
\label{lem:lim}
Let $\gamma\in\R$ and $x:\N\to\R$ with $|\gamma|<1$ and $\{|x_n|:n\in\N\}$ bounded above.
Then $\displaystyle\lim_{n\to\infty}\gamma^n x_n=0$.
\end{lemma}

\section{The Markov assumption in Lean}
\label{sec:markov}

\paragraph{The common property.}
Whatever the application, a first-order Markov assumption says that the probability of a length-$n$ object is a product of $n$ one-step factors, each of which looks back a bounded amount in the right state space.
Two consequences are used throughout this paper: (i) the logarithm of the product is the sum of the logarithms of the factors, so a score is a sum of per-step terms; (ii) the product can be built, summed, maximised or differentiated by induction on the length, which gives linear-time recursions.
What differs from model to model is how the factors are normalized (Section~\ref{sec:tables}).

\subsection{The assumption as a factorization}
Fix a sentence, that is, an observed sequence $xo$, and a finite or countable label type $\Y$.
For label sequences $ys:\N\to\Y$ let $P_t(ys)$ denote the joint probability $\Prb\bigl(x[{:}t{+}1]=xo[{:}t{+}1]\wedge y[{:}t{+}1]=ys[{:}t{+}1]\bigr)$ of the first $t+1$ observations and labels.

\begin{definition}[\texttt{IsHiddenMarkovSeq}]
\label{def:hms}
Let $P:\N\to(\N\to\Y)\to\R$, $\pi:\Y\to\R$, $T:\Y\to\Y\to\R$ and $E:\N\to\Y\to\R$.
Then \texttt{IsHiddenMarkovSeq}$\,P\,\pi\,T\,E$ states that for every $ys$,
\[
P_0(ys)=E_0(ys_0)\,\pi(ys_0),\qquad
P_{t+1}(ys)=P_t(ys)\cdot T(ys_t,ys_{t+1})\,E_{t+1}(ys_{t+1})\quad(t\in\N).
\]
\end{definition}
In the reading we use, $\pi(a)=\Prb(y_0=a)$, $T(a,b)=\Prb(y_{i+1}=b\mid y_i=a)$ and $E_i(b)=\Prb(x_i=xo_i\mid y_i=b)$.
The probabilistic version \texttt{IsHiddenMarkovPr}$\;\pi\;x\;y\;xo$ states the same two equations directly in probability notation: $\Prb(x[{:}1]=xo[{:}1]\wedge y[{:}1]=ys[{:}1])=\Prb(x_0=xo_0\mid y_0=ys_0)\Prb(y_0=ys_0)$ and
\begin{multline*}
\Prb\bigl(x[{:}t{+}2]=xo[{:}t{+}2]\wedge y[{:}t{+}2]=ys[{:}t{+}2]\bigr)\\=\Prb\bigl(x[{:}t{+}1]=\dots\wedge y[{:}t{+}1]=\dots\bigr)\,\Prb(y_{t+1}\mid y_t)\,\Prb(x_{t+1}=xo_{t+1}\mid y_{t+1}).
\end{multline*}
\texttt{IsHiddenMarkovFac} is the same factorization for an abstract prefix family $P$; the theorem \texttt{IsHiddenMarkovPr.toFac} passes from the probabilistic version to it, and \texttt{IsHiddenMarkovSeq.prefix} (also for \texttt{IsHiddenMarkovFac}) proves by induction on $t$ that $P_t$ depends only on the first $t+1$ labels.
These are the definitions and helper theorems that the CRF lemmas below take as hypotheses.

\begin{remark}[Hypothesis versus derivation]
\label{rem:hyp}
The abstract factorization \texttt{IsHiddenMarkovSeq} is \emph{assumed} by the logits lemma of Theorem~\ref{thm:logits0}.
For the probabilistic version \texttt{IsHiddenMarkovPr} the factorization is \emph{derived} from conditional independences.
An observation depends on the past only through the current label, $x_{t+1}\perp(x[{:}t{+}1],y[{:}t{+}1])\mid y_{t+1}$, and a label depends on the past only through its predecessor, $y_{t+1}\perp(x[{:}t{+}1],y[{:}t])\mid y_t$.
Theorem~\ref{thm:markov} shows that these two conditional independences imply the product form, and \lmod{IsHiddenMarkovPr.of.CondIndep.CondIndep}{Random.IsHiddenMarkovPr.of.CondIndep.CondIndep} packages the one-step equations as \texttt{IsHiddenMarkovPr}.
They have to be read as \emph{conditional} independences: the marginal independences (for instance $x_{t+1}\perp(x[{:}t{+}1],y[{:}t{+}1])$ without conditioning on $y_{t+1}$) do not imply the factorization (take $x_0$ constant, $y_0,y_1$ independent fair bits and $x_1=y_0\oplus y_1$, with a little noise to make all probabilities positive).
By contrast, the Markov property of the MDP trajectory law of Section~\ref{sec:mdp} follows from the construction of the law (\texttt{joint\_succ}, \texttt{hist\_step}), not from assumed independences; the treatment of the two sides is asymmetric.
\end{remark}

\subsection{The product form}
\begin{theorem}[\lmod{crf.markov}{Random.Prob.eq.Mul_Prod_MulProbS.of.Ne_0.Eq.Eq.Eq}~\cite{lemma_crf_markov}]
\label{thm:markov}
Let $x_i,y_i$ be measurable discrete random variables (countable ranges, counting reference measures) on a probability space with a standard Borel sample space, and suppose that for every $t$
\[
x_{t+1}\perp\bigl(x[{:}t{+}1],\,y[{:}t{+}1]\bigr)\mid y_{t+1},
\qquad
y_{t+1}\perp\bigl(x[{:}t{+}1],\,y[{:}t]\bigr)\mid y_t .
\]
Then for every $xo$, $ys$ and $t$,
\begin{multline*}
\Prb\bigl(x[{:}t{+}1]=xo[{:}t{+}1]\wedge y[{:}t{+}1]=ys[{:}t{+}1]\bigr)\\
=\Prb(x_0=xo_0\mid y_0=ys_0)\,\Prb(y_0=ys_0)\prod_{i=1}^{t}\Prb(y_i=ys_i\mid y_{i-1}=ys_{i-1})\,\Prb(x_i=xo_i\mid y_i=ys_i).
\end{multline*}
\end{theorem}
The proof is induction on $t$ with \texttt{Finset.prod\_Ico\_succ\_top}.
The one-step identity is the chain rule $\Prb(C\cap F\cap E)=\Prb(C)\,\frac{\Prb(F\cap G)}{\Prb(G)}\,\frac{\Prb(E\cap F)}{\Prb(F)}$ for the past event $C=\{x[{:}t{+}1]=xo[{:}t{+}1],\,y[{:}t{+}1]=ys[{:}t{+}1]\}$, $E=\{x_{t+1}=xo_{t+1}\}$, $F=\{y_{t+1}=ys_{t+1}\}$ and $G=\{y_t=ys_t\}\supseteq C$; it follows by clearing denominators in the two conditional independences.
No non-vanishing hypothesis is needed, because if a conditioning event is null then both sides vanish (the module name keeps the historical \texttt{Ne\_0}).

\subsection{Conditional independence: the standard bricks}
\label{sec:cond}
Three library lemmas formalize the conditional-independence vocabulary in which the Markov assumption is usually stated.
Let $x,y,z$ be measurable random variables on a probability space.
\begin{itemize}[leftmargin=1.6em,itemsep=0.2em]
  \item \lmod{CondIndep.of.Indep_Joint}{Random.CondIndep.of.Indep_Joint}: if $x\perp(y,z)$ then $x\perp y\mid z$.
  \item \lmod{Indep_Joint.of.CondIndep.Indep}{Random.Indep_Joint.of.CondIndep.Indep}: if $z\perp x\mid y$ and $z\perp y$ then $z\perp(x,y)$.
  \item \lmod{All_Imp_EqProbSCond.of.CondIndep}{Random.All_Imp_EqProbSCond.of.CondIndep}: if $x\perp y\mid z$ then, almost everywhere on the reference measures, $\Prb(x\mid y,z)=\Prb(x\mid z)$ whenever $\Prb(y,z)\ne0$; this is the Markov property in conditional form.
\end{itemize}
The chain rule used to factor a joint probability one variable at a time is the discrete lemma \lmod{ProbCond.eq.Mul.ProbCond}{Random.ProbCond.eq.Mul.ProbCond}: for variables with countable ranges and counting reference measures, $\Prb(x\wedge y\mid z)=\Prb(x\mid z)\,\Prb(y\mid x\wedge z)$.
It is two-variable and pointwise; the $n$-step chain rule for the hidden Markov sequence is Theorem~\ref{thm:markov}. It is reduced to the measure-level form \lmod{MulMeasure.eq.MulMeasure.of.CondIndep}{Random.MulMeasure.eq.MulMeasure.of.CondIndep} of conditional independence, $\Prb(x{=}u,y{=}v,z{=}w)\,\Prb(z{=}w)=\Prb(x{=}u,z{=}w)\,\Prb(y{=}v,z{=}w)$ (obtained from the third lemma above, with no non-vanishing hypothesis), and to the identification \lmod{ProbCond.eq.Div.of.Eq_Count.Eq_Count}{Random.ProbCond.eq.Div.of.Eq_Count.Eq_Count} of $\Prb(x\mid y)$ with a ratio of measures.

\section{Forward and backward propagation of the linear-chain CRF}
\label{sec:crf}

This is the first core topic.
A linear-chain CRF assigns to a label sequence $ys\in\Y^{m+1}$ a weight that is a product of one-step factors, and its partition function $Z$ is the sum of these weights over all $k^{m+1}$ label sequences.
The Markov structure makes two sums computable in linear time: the \emph{forward} recursion over prefixes (Section~\ref{sec:forward}) and the \emph{backward} recursion over suffixes (Section~\ref{sec:backward}); a third recursion over the same chain, with $\max$ in place of $\sum$, is Viterbi (Section~\ref{sec:viterbi}).
Sections~\ref{sec:forward} and~\ref{sec:viterbi} are stated for the hidden-Markov-conditional instance in the log domain; Section~\ref{sec:backward} is a purely algebraic statement about arbitrary real weights, to which the hidden Markov factorization is an instance.

\subsection{From the product to the score}
\label{sec:score}
Write $G(a,b)=\log\Prb(y_{i+1}=a\mid y_i=b)$ for the transition score into $a$ from $b$, $e_t(a)=\log\Prb(x_t=xo_t\mid y_t=a)$ for the emission (unary) score, and $s_t(ys)=\log P_t(ys)$ for the log-score of a label prefix.
All probabilities are assumed strictly positive in the following statements.

\begin{theorem}[\lmod{crf.markov.logits}{Tensor.Eq.of.Ne_0.Eq.Eq.Eq.Eq_log.Eq_log.Eq_log}~\cite{lemma_crf_markov_logits}]
\label{thm:logits0}
Under \texttt{IsHiddenMarkovSeq}$\,P\,\pi\,T\,E$ with $P,T,E>0$, $s_t=\log P_t$, $x_t=\log E_t$ and $G(a,b)=\log T(b,a)$, for every $t>0$:
\[
s_t(ys)=G(ys_t,ys_{t-1})+s_{t-1}(ys)+x_t(ys_t).
\]
\end{theorem}

\begin{theorem}[\lmod{crf.logits}{Tensor.Imp.of.Eq}~\cite{lemma_crf_logits}]
\label{thm:logits}
Assume \texttt{IsHiddenMarkovPr}$\;\pi\;x\;y\;xo$, positivity of $\Prb(y_0=a)$, $\Prb(y_{i+1}=a\mid y_i=b)$ and $\Prb(x_t=xo_t\mid y_t=a)$, and the definitions of $s_t$, $e_t$, $G$ above.
Then for every $ys$,
\begin{align*}
s_{t+1}(ys)&=G(ys_{t+1},ys_t)+s_t(ys)+e_{t+1}(ys_{t+1}),\\
s_t(ys)&=\log\Prb(y_0=ys_0)+\sum_{i=1}^{t}G(ys_i,ys_{i-1})+\sum_{i=0}^{t}e_i(ys_i).
\end{align*}
\end{theorem}
The second line is the textbook linear-chain CRF score: a sum of unary (emission) terms and pairwise (transition) terms.
It is exactly the \emph{log of a Markov product is a sum of per-step terms} property of Section~\ref{sec:markov}; the neural CRFs~\cite{huang2015bilstmcrf,lample2016ner,ma2016bilstmcnnscrf} replace $e_t$ by the output of a network and treat $G$ as a free matrix, and the derivations of Su~\cite{su2017crf} start from the same additive form.


\subsection{Forward propagation: log-sum-exp and the forward recursion}
\label{sec:forward}
Two pointwise facts about the softmax are in the library: \lmod{Log.Softmax.eq.Add.LogSumExp}{Tensor.Log.Softmax.eq.Add.LogSumExp}, $\log\frac{e^{x_i}}{\sum_je^{x_j}}=x_i-\log\sum_je^{x_j}$, and \lmod{Le_LogSumExp}{Tensor.Le_LogSumExp}, $x_i\le\log\sum_je^{x_j}$; its Jacobian is \lmod{Grad.Softmax.eq.Mul.Softmax}{Real.Grad.Softmax.eq.Mul.Softmax}: with $\sigma_j=e^{F_j}/\sum_ie^{F_i}$, $\partial\sigma_j=\sigma_j\bigl(\partial F_j-\sum_i\sigma_i\,\partial F_i\bigr)$.

For the forward recursion let $\Y$ be finite with $k$ elements, fix a length $m+1$, and put
\[
z_t(a)=\sum_{ys\in\Y^{t+1},\;ys_t=a}e^{s_t(ys)},\qquad \alpha_t(a)=\log z_t(a).
\]
Thus $z_t(a)$ is a brute-force sum over $k^{t}$ label prefixes, and $Z=\sum_bz_m(b)$ is the sum over all $k^{m+1}$ label paths.

\begin{theorem}[\lmod{y_given_x}{Tensor.Imp_Eq_Add_LogSumExp.EqNegLogProb.of.Ne0ProbJoint.Indep.y_given_x}~\cite{lemma_crf_nll}]
\label{thm:nll}
Under the hypotheses of Theorem~\ref{thm:logits} (with positive prefix probabilities),
\[
\alpha_{t+1}(a)=\log\sum_{b\in\Y}\exp\bigl(\alpha_t(b)+G(a,b)\bigr)+e_{t+1}(a),
\]
and for every label sequence $ys$,
\[
-\log\Prb\bigl(y[{:}m{+}1]=ys[{:}m{+}1]\bigm|x[{:}m{+}1]=xo[{:}m{+}1]\bigr)\;=\;\log\sum_{b\in\Y}e^{\alpha_m(b)}\;-\;s_m(ys).
\]
\end{theorem}
Since $\sum_be^{\alpha_m(b)}=Z$ by definition, the second line is the CRF negative log-likelihood $-\log p(y\mid x)=\log Z(x)-\text{score}(x,y)$, and the first line is the forward recursion that computes $Z$ in $O(mk^2)$ time instead of $k^{m+1}$ terms: Lean checks that the recursion agrees with the brute-force definition of $z_t$ by induction on $t$ (a splitting of the sum over $\Y^{t+1}$ by its last label, \texttt{Fintype.sum\_snoc}).
This is the recursion that computes the normalizer recursively thanks to the Markov assumption, and that Sutton and McCallum~\cite{sutton2012crf} present as the forward algorithm for linear-chain CRFs.
In the probability domain, with $P_t=e^{s_t}$, it reads $z_{t+1}(a)=\bigl(\sum_bz_t(b)\,\Prb(y_{t+1}=a\mid y_t=b)\bigr)\Prb(x_{t+1}=xo_{t+1}\mid y_{t+1}=a)$; this is the form that Theorem~\ref{thm:fb} proves for arbitrary weights, together with its backward counterpart.


\subsection{Backward propagation and the forward--backward identity}
\label{sec:backward}
The forward variable $z_t(a)$ summarizes the prefix $ys_0,\dots,ys_t$ ending in the label $a$.
The \emph{backward} variable $\beta_t(a)$ summarizes the suffix $ys_{t+1},\dots,ys_m$ that may follow the label $a$ at time $t$.
Their product is the total weight of all label sequences through $a$ at time $t$, so $z_t(a)\beta_t(a)/Z$ is the posterior marginal $\Prb(y_t=a\mid x)$ of the hidden-Markov reading.
We state the Lean theorem for arbitrary real weights, because neither positivity nor normalization is needed for the recursions.

\paragraph{Setting.}
Let $\Y$ be a finite type, $\psi_0:\Y\to\R$ and $\psi_t:\Y\to\Y\to\R$ ($t\ge1$) arbitrary real functions, and $P_t:(\N\to\Y)\to\R$ prefix weights with
\begin{equation}
\label{eq:chain}
P_0(ys)=\psi_0(ys_0),\qquad P_{t+1}(ys)=P_t(ys)\,\psi_{t+1}(ys_t,ys_{t+1})\quad(t\in\N).
\end{equation}
The hidden Markov factorization of Definition~\ref{def:hms} is the instance $\psi_0(a)=E_0(a)\pi(a)$ and $\psi_t(a,b)=T(a,b)E_t(b)$ (and \texttt{IsHiddenMarkovFac} is the instance with time-dependent transition probabilities): the two hypotheses of the theorem below are exactly the two fields of these definitions, so the instantiation is the projections \texttt{(h ys).1} and \texttt{(h ys).2 t}.
The same equation~\eqref{eq:chain} with a free positive $\psi_t(y_{t-1},y_t,x)$ at a fixed input $x$ is the unnormalized weight of a linear-chain CRF; the theorem does not define the normalized distribution $\prod_t\psi_t/Z$ as a model.
Put
\[
z_t(a)=\sum_{ys\in\Y^{t+1},\;ys_t=a}P_t(ys)\qquad(\text{for }P_t=e^{s_t}\text{ this is the }z_t\text{ of Section~\ref{sec:forward}}),
\]
fix the length $m+1$, and let the backward variable satisfy
\[
\beta_m(a)=1,\qquad \beta_t(a)=\sum_{c\in\Y}\psi_{t+1}(a,c)\,\beta_{t+1}(c)\quad(t<m).
\]

\begin{theorem}[\lmod{crf.forward_backward}{Tensor.And.of.All_Eq_Sum_Mul.All_Eq_1.All_Eq_Mul.All_Eq.forward_backward}~\cite{lemma_crf_fb}]
\label{thm:fb}
Under the setting above, for every $a\in\Y$:
\begin{enumerate}[leftmargin=1.8em,itemsep=0.2em,label=(\alph*)]
  \item (forward recursion) $z_0(a)=\psi_0(a)$ and $z_{t+1}(a)=\sum_{b\in\Y}z_t(b)\,\psi_{t+1}(b,a)$ for every $t$;
  \item (forward--backward identity) for every $t\le m$,
  \[
  \sum_{ys\in\Y^{m+1},\;ys_t=a}P_m(ys)\;=\;z_t(a)\,\beta_t(a);
  \]
  \item (the partition function does not depend on the cut) for every $t\le m$,
  \[
  \sum_{a\in\Y}z_t(a)\,\beta_t(a)\;=\;\sum_{a\in\Y}z_m(a)\;=\;Z .
  \]
\end{enumerate}
\end{theorem}
Here $Z=\sum_{ys\in\Y^{m+1}}P_m(ys)$ is the partition function; item~(c) is stated in the file with $\sum_az_m(a)$, which equals $Z$ by definition of $z_m$ (grouping the paths by their last label) and is the $Z$ of Theorem~\ref{thm:nll} in the hidden-Markov instance.
The formal statement has no positivity hypothesis and no probability space; the label sequences are the functions $\mathrm{Fin}(m+1)\to\Y$, padded to a function on $\N$ as in the other CRF lemmas, and $\beta$ is any function satisfying the two displayed backward equations.

\paragraph{Proof idea.}
Part~(a) is the splitting of the sum over $\Y^{t+2}$ by its last label (\texttt{Fintype.sum\_snoc}, as in Theorem~\ref{thm:nll}) together with the fact that $P_t$ depends only on the first $t+1$ labels, which is proved by induction on $t$ from~\eqref{eq:chain}.
For part~(b) fix $t$ and $a$ and let $u_j(c)$ be the sum of $P_j(ys)$ over $ys\in\Y^{j+1}$ with $ys_t=a$ and $ys_j=c$ ($j\ge t$).
The same splitting gives $u_{j+1}(c)=\sum_bu_j(b)\psi_{j+1}(b,c)$, and $u_t(c)=z_t(a)$ if $c=a$ and $0$ otherwise.
A downward induction on $t$ shows $\sum_cu_m(c)=\sum_au_t(a)\beta_t(a)$ for every initial vector $u_t$ obeying the recursion; in the inductive step one interchanges two finite sums and applies the backward equation.
Part~(c) is the same induction applied to $u=z$, using $\beta_m=1$.
Everything is a finite-sum manipulation; the cost of the recursions is $O(mk^2)$ for $k=|\Y|$.

\paragraph{Reading.}
In the hidden-Markov reading, $P_m(ys)=\Prb(x[{:}m{+}1]=xo[{:}m{+}1]\wedge y[{:}m{+}1]=ys[{:}m{+}1])$ and $Z=\Prb(x=xo)$, so item~(b) says that $z_t(a)\beta_t(a)$ equals $\Prb(x=xo,\,y_t=a)$, and dividing by $Z$ gives the posterior marginal $\Prb(y_t=a\mid x=xo)$.
The division is the Bayes step that the proof of Theorem~\ref{thm:nll} already carries out for a whole path; the normalized marginal is \emph{not} a separate Lean statement.
Item~(c) is the sanity check that every cut $t$ yields the same $Z$; it is what makes the forward--backward pair a consistent decomposition rather than two unrelated recursions.
Rabiner~\cite{rabiner1989hmm} and Sutton and McCallum~\cite{sutton2012crf} present the forward and backward recursions in the textbook form.
The log-domain version of the backward recursion is the same statement for $\log\beta$ with $\log\sum\exp$; it is not separately formalized.

\subsection{Viterbi: the max-plus recursion}
\label{sec:viterbi}
Replacing sum by maximum and the logarithm by the identity gives the second recursion.
Let $v_t(a)=\max_{ys\in\Y^{t+1},\,ys_t=a}s_t(ys)$.

\begin{theorem}[\lmod{crf.viterbi}{Tensor.And.of.Ne_0.Eq.Eq.Eq}~\cite{lemma_crf_viterbi}]
\label{thm:viterbi}
Under the same hypotheses,
\begin{align*}
v_{t+1}(a)&=e_{t+1}(a)+\max_{b\in\Y}\bigl(v_t(b)+G(a,b)\bigr),\\
\max_{ys\in\Y^{m+1}}\Prb\bigl(x[{:}m{+}1]=xo[{:}m{+}1]&\wedge y[{:}m{+}1]=ys\bigr)=\exp\max_{b\in\Y}v_m(b).
\end{align*}
\end{theorem}
Only the maximal \emph{value} is computed; there is no argmax and no backpointer, so the decoded sequence itself is not formalized.
The forward recursions of Theorems~\ref{thm:nll} and~\ref{thm:viterbi} are two instances of one recursion on the same chain, over the log semiring ($\oplus=$ log-sum-exp, $\otimes=+$) and the max-plus semiring ($\oplus=\max$, $\otimes=+$) respectively; in Lean they are proved separately, and the semiring unification is a remark, not a theorem.
(The quantity $v_t$ is unrelated to the backward variable $\beta_t$ of Section~\ref{sec:backward}.)


\subsection{HMM versus CRF: what the statements say and do not say}
\label{sec:hmmcrf}
\paragraph{HMM.}
A hidden Markov model~\cite{rabiner1989hmm} is \emph{generative} and \emph{locally normalized}:
\[
p(x,y)=\prod_{t}p(y_t\mid y_{t-1})\,p(x_t\mid y_t),\qquad Z=1,
\]
each factor being a probability distribution in its own argument.
\paragraph{Linear-chain CRF.}
A linear-chain CRF~\cite{lafferty2001crf,sutton2012crf} is \emph{discriminative} and \emph{globally normalized}:
\[
p(y\mid x)=\frac1{Z(x)}\prod_{t}\psi_t(y_{t-1},y_t,x),\qquad Z(x)=\sum_{y'\in\Y^{n}}\prod_t\psi_t(y'_{t-1},y'_t,x),
\]
with arbitrary non-negative potentials $\psi_t$ that need not be probabilities, and one partition function $Z(x)$ that is a sum over $k^n$ label paths.
The Markov assumption is what makes $Z(x)$ computable by the forward recursion of Theorem~\ref{thm:nll}, or equivalently by the forward--backward pair of Theorem~\ref{thm:fb}.
In terms of the score recursion of Theorem~\ref{thm:logits}, a CRF is the HMM score recursion with the local-normalization constraint removed: the potentials $\psi_t=e^{G+e_t}$ are free, and what remains of the probability is $e^{\text{score}}/Z(x)$.
\paragraph{What Lean covers.}
Theorems~\ref{thm:logits}--\ref{thm:nll} and~\ref{thm:viterbi} are the HMM-\emph{conditional} instance: $\psi_t(y_{t-1},y_t,x)=\Prb(y_t\mid y_{t-1})\Prb(x_t=xo_t\mid y_t)$ with positive probabilities, for which $Z(x)$ is the evidence $\Prb(x=xo)$.
Theorem~\ref{thm:fb} covers the forward and backward recursions for arbitrary real weights $\psi_t$, hence the algebra of a general-potential CRF at a fixed input; but a general-potential CRF (arbitrary $\psi_t$, a trainable transition matrix that is not a conditional distribution, hence a genuinely global $Z$) is \emph{not} defined as a model, and the quantities that need that model are not formalized.
The Lean conclusion $-\log p=\log Z-\text{score}$ is therefore the statement that makes the CRF loss precise in the case that is provable from the factorization hypothesis, not a proof about CRFs in general.

\section{The Bellman equation}
\label{sec:bellman}

This is the second core topic.
The CRF lemmas fix an observed sentence and sum over labels; here the randomness is a whole trajectory, and the Markov structure is that of a decision process: the next state and reward depend on the past only through the current state and action.
The Bellman equation says that the value of a state is the expected immediate reward plus the discounted value of the next state, a one-step recursion that replaces an infinite-horizon expectation.

\subsection{Model and trajectory measure}
\label{sec:mdp}
\begin{definition}[Model]
\label{def:model}
Let $S$ and $A$ be finite types with the discrete $\sigma$-algebra and $\Theta$ a real normed space.
A \emph{policy} is $\pi:\Theta\to S\to A\to\R$, written $\pi_\theta(u\mid x)$, with $\pi_\theta(u\mid x)\ge0$ and $\sum_u\pi_\theta(u\mid x)=1$ (local normalization: $Z=1$ at every step).
An \emph{environment} consists of an initial probability measure $\iota$ on $S$, a Markov kernel $T$ from $S\times A$ to $S$, a Markov kernel $R$ from $S\times A$ to $\R$ with support in $[-R_{\max},R_{\max}]$.
\end{definition}
A trajectory is a sequence $\omega\in(S\times A\times\R)^{\N}$ of stages with coordinates $\sv_t,\av_t,\rw_t$.
One stage from $x$ is drawn by $\kappa_\theta(x)=\delta_x\otimes(\pi_\theta(\cdot\mid x)\otimes R)$ and the stage chain has kernel $K_\theta(x,u,\rho)=(\kappa_\theta\circ T)(x,u)$.
The trajectory measure $\Prb_\theta=\mathtt{Kernel.trajMeasure}\;\mu_0\;(K_\theta)_n$ is the Ionescu-Tulcea extension~\cite{kallenberg2002foundations} provided by mathlib, with $n$-th kernel reading only the last stage of the history.
Informally (the product formula is not a separate Lean statement), the finite-dimensional marginals have the Markov-product shape $\iota(s_0)\prod_t\pi_\theta(a_t\mid s_t)\,T(s_{t+1}\mid s_t,a_t)\,R(\cdot\mid s_t,a_t)$.
Its logarithm is a sum of per-step terms, and since the transition and reward factors do not depend on $\theta$, \emph{the score of a trajectory is the sum of the scores of its actions}, $\sum_t\grad\log\pi_\theta(a_t\mid s_t)$.
This is the MDP counterpart of the CRF score of Theorem~\ref{thm:logits}; Lean uses it only through the expectation lemma \texttt{E\_score} below.

\paragraph{The Markov property is a theorem here.}
In contrast to Remark~\ref{rem:hyp}, where the factorization follows from assumed conditional independences, the Markov property of $\Prb_\theta$ is proved from the construction.
The model file contains \texttt{joint\_succ}: the law of $(\omega_t,\omega_{t+1})$ is the law of $\omega_t$ composed with the kernel $K_\theta$, and \texttt{hist\_step}: for a bounded measurable $\varphi$,
\[
\E\bigl[\varphi(h_n,\omega_{n+1})\bigr]=\int\!\Bigl(\int\varphi(h,z)\,dK_\theta(\mathrm{last}\,h)(z)\Bigr)\,d\,\mathrm{law}(h_n)(h),
\]
where $h_n=(\omega_0,\dots,\omega_n)$ is the history and $\mathrm{last}\,h$ is its last stage, together with iterated versions (\texttt{hist\_mul}, \texttt{hist\_iter}) used below.
The joint law of state and action factors as $\Prb(\sv_t=x,\av_t=u)=\Prb(\sv_t=x)\,\pi_\theta(u\mid x)$ (\lmod{ProbJoint.eq.Mul_Prob_Pr}{Random.ProbJoint.eq.Mul_Prob_Pr}).


\subsection{Value functions and the Bellman equations}
\label{sec:bellmaneq}
For $\gamma\in[0,1)$, the discounted return from $t$ is $G_t=\tsum_k\gamma^k\rw_{t+k}$, and
$V_t^\theta(s_t)=\E_\theta[G_t\mid\sv_t]$, $Q_t^\theta(s_t,a_t)=\E_\theta[G_t\mid\sv_t,\av_t]$.
The objective is $J(\theta)=\E_\theta[\tsum_t\gamma^t\rw_t]=\tsum_t\gamma^t\E_\theta[\rw_t]$.
The model file shows that on reachable states $V_t^\theta$ equals a time-free closed form, that $|V_t|,|Q_t|\le(1-\gamma)^{-1}R_{\max}$, and that $J$ is Fr\'echet differentiable for a differentiable policy with uniformly bounded gradient.
We assume throughout:
(D) $\gamma\in[0,1)$; (P1) $\theta\mapsto\pi_\theta(u\mid x)$ is differentiable; (P2) $\sup_{\theta,x,u}\|\grad\pi_\theta(u\mid x)\|<\infty$.
Bounds on $V_t$ and $\grad V_t$ are derived, not assumed.

\begin{theorem}[Bellman equations~\cite{lemma_bellman}]
\label{thm:bellman}
Assume (D) and let $V,Q$ satisfy $V_t(s_t)=\E_\theta[G_t\mid\sv_t]$, $Q_t(s_t,a_t)=\E_\theta[G_t\mid\sv_t,\av_t]$.
Then
\begin{align*}
V_t(s_t)&=\E_\theta\bigl[Q_t(s_t,\av_t)\bigm|\sv_t\bigr],\\
V_t(s_t)&=\E_\theta\bigl[\rw_t+\gamma\,V_{t+1}(\ra{s}_{t+1})\bigm|\sv_t\bigr],\\
Q_t(s_t,a_t)&=\E_\theta\bigl[\rw_t+\gamma\,V_{t+1}(\ra{s}_{t+1})\bigm|\sv_t,\av_t\bigr].
\end{align*}
\end{theorem}
In the Lean file the three identities are stated for every $t$ and every pair $(x,u)$ with no reachability hypothesis: at a null conditioning event both sides are $0$ (Section~\ref{sec:notation}).
The proof combines the tower property of conditional expectation with the Markov property of the trajectory law of the previous subsection.
Read from $t+1$ to $t$, the second identity is a \emph{backward} recursion for the value function, the same direction as the recursion for $\beta_t$ in Theorem~\ref{thm:fb}; Section~\ref{sec:analogy} compares the two without claiming a theorem.
The gradient of this recursion is the first step of the policy gradient theorem (Theorem~\ref{thm:rec}).


\subsection{Forward, backward and Bellman recursions side by side (analogy)}
\label{sec:analogy}
This subsection is an analogy, not a theorem: no Lean statement relates the CRF recursions to the Bellman recursion.
Table~\ref{tab:rec} lists the four one-step recursions of this paper, the aggregation each performs over the Markov chain, the direction in which it runs, and the Lean statement that proves it.

\begin{table}[ht]
\centering
\footnotesize
\renewcommand{\arraystretch}{1.2}
\begin{tabular}{@{}>{\raggedright\arraybackslash}p{0.12\linewidth}>{\raggedright\arraybackslash}p{0.34\linewidth}>{\raggedright\arraybackslash}p{0.10\linewidth}>{\raggedright\arraybackslash}p{0.13\linewidth}>{\raggedright\arraybackslash}p{0.21\linewidth}@{}}
\toprule
Recursion & One step & Direction & Aggregation & Lean \\
\midrule
CRF forward & $z_{t+1}(a)=\sum_bz_t(b)\,\psi_{t+1}(b,a)$ & forward in $t$ & sum over prefixes (sum-product) & Theorem~\ref{thm:fb}(a); log form Theorem~\ref{thm:nll} \\
CRF backward & $\beta_t(a)=\sum_c\psi_{t+1}(a,c)\,\beta_{t+1}(c)$ & backward in $t$ & sum over suffixes (sum-product) & Theorem~\ref{thm:fb}(b),(c) \\
Viterbi & $v_{t+1}(a)=e_{t+1}(a)+\max_b\bigl(v_t(b)+G(a,b)\bigr)$ & forward in $t$ & max over prefixes (max-sum) & Theorem~\ref{thm:viterbi} (value only) \\
Bellman & $V_t(s_t)=\E_\theta[\rw_t+\gamma V_{t+1}(\ra{s}_{t+1})\mid\sv_t]$ & backward in $t$ & expectation over the next state and reward & Theorem~\ref{thm:bellman} \\
\bottomrule
\end{tabular}
\caption{One-step recursions on a Markov chain. The first three are dynamic programs over label paths; the Bellman recursion is a dynamic program over trajectories.}
\label{tab:rec}
\end{table}

Formally, if one drops the reward in the Bellman recursion and imposes a terminal value, it becomes $V_t(s)=\sum_cK(s,c)\,V_{t+1}(c)$ for the one-step kernel $K$ of the state chain, which has exactly the shape of the backward recursion $\beta_t(a)=\sum_c\psi_{t+1}(a,c)\beta_{t+1}(c)$ with $\psi\leftrightarrow K$ and $\beta_m\leftrightarrow V_m$.
The differences are real: $K$ is a probability kernel ($\sum_cK(s,c)=1$) whereas $\psi$ is an unnormalized weight, the Bellman recursion carries a reward source term and a discount, and the Lean value function is an infinite-horizon conditional expectation, not a finite suffix sum.
What the recursions share is the mechanism: a Markov structure replaces a global sum by a one-step recursion, forward when the quantity is built from the past, backward when it is built from the future.

\section{The policy gradient theorem: REINFORCE and GAE}
\label{sec:pg}

This is the third core topic.
The policy gradient theorem expresses the gradient of the expected discounted return as an expectation of the return times the score $\grad\log\pi_\theta(\av_t\mid\sv_t)$ of the action actually taken.
Lean proves it by induction on the horizon from the Bellman recursion, in three forms: a truncated form with an explicit remainder (Theorem~\ref{thm:trunc}), an action-value form, and the REINFORCE form (Theorem~\ref{thm:pgt}).
Because local normalization makes the score zero-mean (Lemma~\ref{lem:zero}), the return may be replaced by an advantage, which gives the unbiased advantage estimate and GAE (Section~\ref{sec:gae}).

\subsection{The policy gradient theorem by induction}
\label{sec:pgind}
\begin{theorem}[Recursion~\cite{lemma_recursion}, unrolling~\cite{lemma_induct}]
\label{thm:rec}
Assume (D), (P1), (P2) and let $s_0$ be reachable.
Then $\grad V_t(s_t)=\sum_uQ_t(s_t,u)\grad\pi_\theta(u\mid s_t)+\gamma\sum_y\Prb_\theta(\sv_{t+1}=y\mid\sv_t)\grad V_{t+1}(y)$ at reachable $s_t$, and for every $n$,
\[
\grad V_0(s_0)=\sum_{t<n}\gamma^t\sum_y\Prb_\theta(\sv_t=y\mid\sv_0)\sum_uQ_t(y,u)\grad\pi_\theta(u\mid y)+\gamma^n\sum_y\Prb_\theta(\sv_n=y\mid\sv_0)\grad V_n(y).
\]
\end{theorem}
The second identity is proved by induction on $n$, expanding the remainder with the first and collapsing the double sum by the Chapman--Kolmogorov identity.

\begin{theorem}[Truncated policy gradient~\cite{lemma_pgn}]
\label{thm:trunc}
Assume (D), (P1), (P2).
For every $n\in\N$,
\[
\tsum_t\gamma^t\grad\E_\theta[\rw_t]=\E_\theta\Bigl[\sum_{t<n}\gamma^tQ_t(\ra{s}_t,\ra{a}_t)\,\grad\log\pi_\theta(\av_t\mid\sv_t)\Bigr]+\gamma^n\,\E_\theta\bigl[\grad V_n(\ra{s}_n)\bigr].
\]
\end{theorem}
The proof evaluates $\E_\theta[g(\sv_t,\av_t)\grad\log\pi_\theta(\av_t\mid\sv_t)]=\sum_y\Prb_\theta(\sv_t=y)\sum_ug(y,u)\grad\pi_\theta(u\mid y)$, the log-derivative trick $\pi\grad\log\pi=\grad\pi$ against the law of $(\sv_t,\av_t)$.
Despite the finite sum, the theorem concerns the discounted infinite-horizon objective: it is exact for every truncation level, with an explicit remainder.
Letting $n\to\infty$, with the remainder going to $0$ by Lemma~\ref{lem:lim} and the bound on $\grad V_n$ over reachable states, gives the action-value form~\cite{lemma_pgq}
$\tsum_t\gamma^t\grad\E_\theta[\rw_t]=\tsum_t\gamma^t\E_\theta[Q_t(\ra{s}_t,\ra{a}_t)\grad\log\pi_\theta(\av_t\mid\sv_t)]$, and with the tower property~\cite{lemma_tower}, $\E_\theta[G_t\grad\log\pi_\theta]=\E_\theta[Q_t\grad\log\pi_\theta]$, the REINFORCE form.

\begin{theorem}[Policy gradient, REINFORCE form~\cite{lemma_pgt}]
\label{thm:pgt}
Let $\Theta$ be a real Hilbert space, $S$ and $A$ carry counting measure, and assume (D), (P1), (P2).
Then
\[
\grad\E_\theta\Bigl[\tsum_t\gamma^t\rw_t\Bigr]=\E_\theta\Bigl[\tsum_t\gamma^t\,G_t\,\grad\log\pi_\theta(\av_t\mid\sv_t)\Bigr].
\]
\end{theorem}
Here the series sit inside the expectation; this is possible because rewards are bounded, and the shared lemmas \lmod{Log.Pr.of.Bounded}{Tensor.Eq.Expect.Sum.Grad.Log.Pr.of.Bounded} and \lmod{Bounded.Discount}{Tensor.Eq.Expect.Sum.Grad.Log.Pr.of.Bounded.Discount} exchange series and expectations.


\subsection{Normalization gives a zero-mean score}
\label{sec:zeroscore}
The property that links the local normalization of Definition~\ref{def:model} to variance reduction is the following.

\begin{lemma}[\lmod{Expect_ConditionedGrad_LogProb.eq.Zero}{Random.Expect_ConditionedGrad_LogProb.eq.Zero}~\cite{lemma_zero}]
\label{lem:zero}
Let $p:\Theta\to A\to\R$ with $\sum_ap_{\theta'}(a)=1$ for all $\theta'$, $p_\theta(a)>0$ and $\theta'\mapsto p_{\theta'}(a)$ differentiable at $\theta$.
Then $\displaystyle\E_{a\sim p_\theta}\bigl[\grad\log p_\theta(a)\bigr]=\sum_ap_\theta(a)\,\grad\log p_\theta(a)=0$.
\end{lemma}
The proof is one line: $p\grad\log p=\grad p$, and $\grad\sum_ap=\grad1=0$.
Its conditional form for the policy of the trajectory model is \lmod{LogProb.eq.Zero.policy}{Random.Expect_ConditionedGrad_LogProb.eq.Zero.policy}~\cite{lemma_score}: $\E_\theta[\grad\log\pi_\theta(\av_t\mid\sv_t)\mid\sv_t=x]=0$ for a positive differentiable policy.
Inside the model file the same fact is used without positivity (terms with $\pi_\theta(u\mid x)=0$ vanish): for every $h:S\to\R$, $\E_\theta[h(\sv_t)\,\grad\log\pi_\theta(\av_t\mid\sv_t)]=0$ (\texttt{E\_h\_score}).
This is why a baseline depending only on the state leaves the policy gradient unchanged, the step used in Section~\ref{sec:gae}.
For a globally normalized CRF the analogue is different: $\E_{p(y\mid x)}[\grad\log p(y\mid x)]=0$ holds as well, but $\grad\log p=\grad\mathrm{score}-\E[\grad\mathrm{score}]$ contains the partition function; the identity $\grad\log Z=\E[\grad\mathrm{score}]$ is not formalized.


\subsection{Unbiased advantage and generalized advantage estimation}
\label{sec:gae}

Let $\delta_j=\rw_j+\gamma V_{j+1}(\sv_{j+1})-V_j(\sv_j)$ be the temporal-difference residual of a value function $V$.

\begin{theorem}[Unbiased advantage estimate~\cite{lemma_uae}]
\label{thm:main}
Let $\Theta$ be a real Hilbert space, $S$ and $A$ carry counting measure, assume (D), (P1), (P2), and let $V^\theta_t(s_t)=\E_\theta[G_t\mid\sv_t]$ for all $\theta,t$ and all reachable $s_t$.
With $\hat A_t=\tsum_k\gamma^k\delta_{t+k}$,
\[
\grad\E_\theta\Bigl[\tsum_t\gamma^t\rw_t\Bigr]=\E_\theta\Bigl[\tsum_t\gamma^t\,\hat A_t\,\grad\log\pi_\theta(\av_t\mid\sv_t)\Bigr].
\]
\end{theorem}
For bounded $V$ the series telescopes to $\hat A_t=G_t-V_t(\sv_t)$, so Theorem~\ref{thm:main} reduces to Theorem~\ref{thm:pgt} plus the vanishing of the baseline term $\E_\theta[V_t(\sv_t)\grad\log\pi_\theta]=0$ of Section~\ref{sec:zeroscore}.

\begin{theorem}[Generalized advantage estimate~\cite{lemma_gae}]
\label{thm:gae}
Under the hypotheses of Theorem~\ref{thm:main}, let $\lambda\in[0,1]$ and $\hat A^{\lambda}_t=\tsum_k(\gamma\lambda)^k\delta_{t+k}$.
Then
\[
\grad\E_\theta\Bigl[\tsum_t\gamma^t\rw_t\Bigr]=\E_\theta\Bigl[\tsum_t\gamma^t\,\hat A^{\lambda}_t\,\grad\log\pi_\theta(\av_t\mid\sv_t)\Bigr].
\]
\end{theorem}
For $\lambda\in[0,1)$ the weighted-average form $\hat A^{\lambda}_t=(1-\lambda)\sum_{k\ge0}\lambda^k\sum_{i\le k}\gamma^i\delta_{t+i}$ of Schulman et al.~\cite{schulman2016gae} gives the same identity~\cite{lemma_gae_avg}, via the series identity~\cite{lemma_gae_series}.
The reason is that, for the exact value function, $\E[\delta_{n+1}\psi(\sv_t,\av_t)]=0$ for $t\le n$ (Markov property and Bellman equation), so only the first residual survives against the score.
The statements concern the exact value function: with an approximate critic the terms $k\ge1$ need not have zero mean and $\lambda<1$ is biased~\cite{schulman2016gae}; that is not formalized.

\section{PPO: what Lean covers and what it does not}
\label{sec:ppo}

\paragraph{MDP versus PPO.}
The MDP is a \emph{model}; PPO is an \emph{algorithm} defined on top of it.
Everything in Sections~\ref{sec:bellman} and~\ref{sec:pg} is about the model: the trajectory measure and its Markov property, the Bellman equations, the policy gradient theorem, the advantage and GAE.
PPO~\cite{schulman2017ppo} adds an importance ratio $\rho_t(\theta)=\pi_\theta(a_t\mid s_t)/\pi_{\theta_{\rm old}}(a_t\mid s_t)$, a clipped surrogate
\[
L^{\rm CLIP}(\theta)=\hat{\E}_t\bigl[\min\bigl(\rho_t(\theta)\hat A_t,\;\mathrm{clip}(\rho_t(\theta),1-\epsilon,1+\epsilon)\hat A_t\bigr)\bigr],
\]
several epochs of minibatch stochastic updates on the same rollouts, a learned value network trained with a value loss (often clipped), an entropy bonus, and in RLHF~\cite{ziegler2019finetuning,ouyang2022instructgpt} a per-token KL penalty to a reference policy.
None of these is formalized.

\paragraph{What the theorems give for PPO.}
(i) The advantage estimator of PPO is a truncated GAE (its Eqs.~(10)--(11)); Theorem~\ref{thm:gae} justifies, for the exact value function, the \emph{untruncated} estimator as an unbiased weight of the score inside the policy gradient.
(ii) At $\theta=\theta_{\rm old}$ the ratio is $1$ and $\grad\rho_t=\grad\log\pi_\theta(a_t\mid s_t)$, so the gradient of the unclipped surrogate is the integrand of Theorem~\ref{thm:gae}; this one-line calculus remark is not a Lean statement.
(iii) The language-model case is an MDP with state $(x,y_{<t})$, action the next token, deterministic transition and the reward at the last token; this instance is not formalized either (Section~\ref{sec:gen}).

\paragraph{What is not covered.}
There is no clipped surrogate or importance ratio, no TRPO-style KL constraint or KL penalty (the library has $\mathrm{KL}\ge0$, \lmod{Random.KL.ge.Zero}{Random.KL.ge.Zero}, which does not enter any PPO statement), no value clipping, no learned critic or its bias, no truncated-horizon GAE, no reward model, no minibatch sampling, and no monotone-improvement or convergence claim.
Lean justifies the (truncated) GAE advantage estimator \emph{inside the policy gradient}; it does not justify the clipped objective.

\section{Beyond the core: labeling versus generation (conceptual)}
\label{sec:gen}

This section is a comparison, not a theorem, and it is not part of the three core topics.
The library has no sequence-to-sequence, teacher-forcing, exposure-bias or label-bias lemma; what it has are the two building blocks that the comparison combines.

\paragraph{Two normalizations of the same numerator.}
Let $\mathrm{score}(x,y)=\sum_tf_t(y_{<t},y_t;x)$ be an additive score.
A \emph{globally normalized} model sets $p(y\mid x)=e^{\mathrm{score}(x,y)}/Z(x)$ with $Z(x)=\sum_{y'}e^{\mathrm{score}(x,y')}$, so that $-\log p=\log Z(x)-\mathrm{score}(x,y)$; for a linear chain, Theorem~\ref{thm:nll} computes $Z$ exactly.
A \emph{locally normalized} model normalizes every step by its own softmax, $p(y\mid x)=\prod_tp(y_t\mid x,y_{<t})$ with $p(y_t\mid\cdot)=e^{f_t}/\sum_ve^{f_t(v)}$, so that
\[
-\log p(y\mid x)=\sum_t\Bigl(\log\sum_ve^{f_t(v)}-f_t(y_t)\Bigr)=\sum_t\bigl(\log Z_t-f_t(y_t)\bigr).
\]
Each summand is exactly the Lean identity \lmod{Log.Softmax.eq.Add.LogSumExp}{Tensor.Log.Softmax.eq.Add.LogSumExp} of Section~\ref{sec:crf}; the product-to-sum step is the same \emph{log of a Markov product is a sum} property.
The two models have the same numerator and differ only in the denominator: one $Z(x)$ against a product of per-step normalizers.
The MEMM~\cite{mccallum2000memm} is the locally normalized sibling of the CRF, and its label-bias problem (a state with few outgoing transitions passes all its mass on regardless of the observation) is the counterpart of the exposure bias of sequence-to-sequence models.
Andor et al.~\cite{andor2016global} make the same case for globally normalized neural parsers.

\paragraph{What changes from labeling to generation.}
\begin{enumerate}[leftmargin=1.6em,itemsep=0.2em]
  \item \emph{Alignment and observation.} In labeling, $|y|=|x|$ and $x$ is observed in full at every position; in generation the output length is variable, there is no alignment, and the generated prefix $y_{<t}$ is part of the conditioning of step $t$.
  \item \emph{The Markov structure is relocated.} A linear-chain CRF is first-order in the \emph{labels} and arbitrary in $x$; an autoregressive model $p(y_t\mid x,y_{<t})$ conditions on the \emph{whole prefix}, so it is Markov only in the enlarged state $(x,y_{<t})$.
  In that state space the transition is deterministic (append the chosen token), and the Markov property is trivial.
  \item \emph{Normalizer.} $Z(x)$ over $k^n$ paths is computable by dynamic programming in $O(nk^2)$ because of first-order dependence; for a prefix-conditioned model the global sum over $V^n$ sequences has no such recursion, which is why one normalizes locally.
  \item \emph{Inference.} Viterbi is exact for the chain; greedy, beam and sampling are approximate for generation.
  \item \emph{Training.} Labeling minimizes the supervised $\log Z-\mathrm{score}$. Generation is trained by teacher-forced cross-entropy, which needs the true prefix, or by expected reward, which is where policy gradients enter (Section~\ref{sec:pg}).
\end{enumerate}

\paragraph{A beam-search example.}
Consider a two-step example with $p(a)=0.6$, $p(c)=0.4$, $p(b\mid a)=0.55$, $p(d\mid a)=0.45$, $p(b\mid c)=0.1$, $p(d\mid c)=0.9$.
Greedy decoding (beam size $1$) returns $(a,b)$ with probability $0.33$; beam size $2$ returns $(c,d)$ with probability $0.36$, the global maximum among the four sequences ($0.33,\,0.27,\,0.04,\,0.36$).
A locally normalized model is committed to its per-step decisions.
This example is not checked in Lean.

\section{Comparison tables}
\label{sec:tables}

\subsection{The Markov assumption across models}
Table~\ref{tab:markov} compares what is factorized, what each factor conditions on, how it is normalized, and what is in Lean.
The rows are the models of this paper; \emph{conditioned on} is the information a factor may use.
In every row the probability of a length-$n$ object is a product of one-step factors, which gives (i) a sum of local terms in the log, (ii) a linear-time recursion (forward, backward, Viterbi, Bellman, value iteration) and (iii) induction on $t$ as the proof technique (induction on $t$ for the CRF lemmas; \texttt{hist\_step} and the unrolling for the MDP).

\begin{table}[ht]
\centering
\footnotesize
\renewcommand{\arraystretch}{1.15}
\begin{tabular}{@{}>{\raggedright\arraybackslash}p{0.115\linewidth}>{\raggedright\arraybackslash}p{0.20\linewidth}>{\raggedright\arraybackslash}p{0.16\linewidth}>{\raggedright\arraybackslash}p{0.12\linewidth}>{\raggedright\arraybackslash}p{0.25\linewidth}@{}}
\toprule
Model & Factorized & Conditioned on & Normalization & Lean \\
\midrule
HMM & $p(x,y)=\prod_tp(y_t\mid y_{t-1})\,p(x_t\mid y_t)$ (generative) & previous label; emission on current label only & local, $Z=1$ & \texttt{IsHiddenMarkovSeq}/\texttt{Pr}, \texttt{crf.markov} (from conditional independences) \\
Linear-chain CRF & $p(y\mid x)=\frac1{Z(x)}\prod_t\psi_t(y_{t-1},y_t,x)$ (discriminative) & adjacent labels; $x$ arbitrary, fully observed & global $Z(x)$ over $k^n$ paths & HMM-conditional instance: \texttt{crf.logits}, $\log Z-$score, \texttt{crf.viterbi}; forward and backward recursions for arbitrary weights (\texttt{crf.forward\_backward}); general $\psi_t$ as a model, normalized marginals and $\nabla\log Z$ not formalized \\
MEMM & $\prod_tp(y_t\mid y_{t-1},x)$ & previous label and $x$ & local, per-step softmax & not formalized \\
MDP / policy gradient & $\iota(s_0)\prod_t\pi_\theta(a_t\mid s_t)T(s_{t+1}\mid s_t,a_t)$ & next state and reward depend on $(s_t,a_t)$ only; policy on $s_t$ only & local, $Z=1$ & trajectory measure, \texttt{joint\_succ}, \texttt{hist\_step}, Bellman, policy gradient, advantage, GAE \\
PPO (algorithm) & same trajectory law; objective $L^{\rm CLIP}$ with ratio $\rho_t$ and learned critic & as MDP & local, $Z=1$ & \textbf{none}: clipped surrogate, ratio, KL, value clipping, learned critic not formalized; only the (truncated) GAE advantage inside the policy gradient is justified (exact $V$) \\
Autoregressive LM & $\prod_tp(y_t\mid x,y_{<t})$ & full prefix & local, per-step softmax & not formalized (only log-softmax and the two-variable chain rule); as an MDP with state $(x,y_{<t})$ the Markov property is trivial \\
\bottomrule
\end{tabular}
\caption{The Markov assumption in the models of this paper. HMM versus CRF differ in normalization (local versus global) and in what is modelled ($p(x,y)$ versus $p(y\mid x)$); MDP versus PPO differ in kind (model versus algorithm).}
\label{tab:markov}
\end{table}

\subsection{Sequence labeling versus sequence-to-sequence generation}
Table~\ref{tab:labgen} summarizes Section~\ref{sec:gen}; it is a conceptual comparison.
Only the entries marked ``Lean'' are theorems; everything else is a conceptual comparison.

\begin{table}[ht]
\centering
\footnotesize
\renewcommand{\arraystretch}{1.15}
\begin{tabular}{@{}>{\raggedright\arraybackslash}p{0.14\linewidth}>{\raggedright\arraybackslash}p{0.37\linewidth}>{\raggedright\arraybackslash}p{0.37\linewidth}@{}}
\toprule
 & Sequence labeling (CRF) & Seq2seq generation (LM / PPO) \\
\midrule
Input / alignment & aligned, $|y|=|x|$, $x$ fully observed & unaligned, variable length; the generated prefix is part of the conditioning \\
Normalization & global partition function over $k^n$ label sequences & locally normalized token softmax \\
Dependence & adjacent labels (first order) and all of $x$ & full prefix $y_{<t}$ (and $x$); Markov only in the state $(x,y_{<t})$ \\
Training & supervised likelihood, $\text{NLL}=\log Z-\text{score}$ (Lean, HMM-conditional case) & teacher-forced cross-entropy, or expected reward (policy gradient / PPO) \\
Inference & exact dynamic programming: forward, backward, Viterbi (Lean: forward and backward sums; Viterbi values only) & approximate: greedy, beam, sampling \\
Failure mode & label bias if locally normalized (MEMM); invalid tag sequences with independent softmax & exposure bias; beam-size anomaly (the example of Section~\ref{sec:gen}) \\
Gradient & $\grad\text{NLL}=\E_{\rm model}[\text{features}]-\text{empirical}$ (not in Lean) & $\grad\E[R]=\E[\sum_t\gamma^t\hat A_t\grad\log\pi_\theta(a_t\mid s_t)]$ (Lean, finite $S,A$, exact $V$); $\E[\grad\log\pi]=0$ (Lean) \\
\bottomrule
\end{tabular}
\caption{Sequence labeling versus sequence-to-sequence generation.}
\label{tab:labgen}
\end{table}

\section{Formalization artifact}
\label{sec:artifact}

\paragraph{Modules and version.}
Each theorem links to its module in~\cite{lemma_repo} (prefix \texttt{Lemma.} omitted).
The statements of this paper are those of the working tree following commit \texttt{5a58c908e}: the CRF lemmas (Section~\ref{sec:crf}), including the new forward--backward lemma of Theorem~\ref{thm:fb}, the file with \texttt{IsHiddenMarkovSeq}, the advantage and GAE files and their helper lemmas were uncommitted when this was written.
The temporal-difference lemmas (\texttt{delta\_ae\_bdd}, \texttt{E\_delta\_psi}, \texttt{E\_sum\_delta}) are in the advantage file of the trajectory model; the bounds \texttt{V\_bdd} and \texttt{gradV\_bdd} used in place of hypotheses on $V$ and $\grad V$ are in the differentiability file.
The project uses Lean \texttt{v4.33.1} with mathlib (revision \texttt{0df444a}).

\paragraph{Checking.}
None of the cited modules contains \texttt{sorry}, \texttt{admit} or a new \texttt{axiom}.
\texttt{\#print axioms} reports for the declarations of Theorems~\ref{thm:markov}--\ref{thm:viterbi}, for the log-softmax, conditional-independence, chain-rule and zero-score lemmas of Sections~\ref{sec:markov}--\ref{sec:zeroscore}, and for Theorems~\ref{thm:bellman}--\ref{thm:gae} and the helper lemmas, the axioms \texttt{propext}, \texttt{Classical.choice}, \texttt{Quot.sound} and nothing else.

\section{Discussion, limitations and conclusion}
\label{sec:discussion}

\paragraph{Regularity.}
The reinforcement-side theorems use finite $S$ and $A$ with counting measure, bounded rewards, $\gamma<1$, differentiability and a uniform gradient bound (P1), (P2).
The bounds on $V_t$ and $\grad V_t$ are derived from these, not assumed.
The labeling-side theorems assume strictly positive probabilities so that logarithms are defined.

\paragraph{What is not formalized.}
We list the gaps that matter for the story of this paper.
\begin{itemize}[leftmargin=1.6em,itemsep=0.1em]
  \item \emph{General CRF.} The log-domain lemmas cover only the hidden-Markov-conditional instance; a CRF with arbitrary positive potentials $\psi_t(y_{t-1},y_t,x)$ and its normalization $Z(x)$ is not defined as a model. The Markov structure enters the downstream lemmas as a factorization hypothesis (\texttt{IsHiddenMarkovSeq}, \texttt{IsHiddenMarkovPr}, or the chain~\eqref{eq:chain} of Theorem~\ref{thm:fb}); \texttt{IsHiddenMarkovPr} is derived from conditional independences (Theorem~\ref{thm:markov}), but no CRF construction yields it.
  \item \emph{Inference.} The forward recursion, the backward recursion and the forward--backward identity (Theorem~\ref{thm:fb}) and the Viterbi \emph{value} are covered. Not covered: the normalized posterior and the pairwise marginals, the gradient $\grad\log Z=\E[\grad\mathrm{score}]$, the log-domain backward recursion, and an argmax Viterbi path with back-pointers.
  \item \emph{Generation.} There is no sequence-to-sequence model, no teacher-forcing lemma, no statement about label bias or exposure bias, and no language-model-as-MDP instance. Section~\ref{sec:gen} is a comparison, not a theorem.
  \item \emph{Analogy.} The correspondence between the forward, backward and Bellman recursions (Section~\ref{sec:analogy}) is not a Lean statement.
  \item \emph{PPO.} No clipped objective, importance ratio, KL or value clipping, learned critic, truncated GAE, reward model, monotone improvement or convergence (Section~\ref{sec:ppo}).
  \item \emph{GAE.} Only the exact value function; $\lambda<1$ with an approximate critic introduces bias that is not analysed.
  \item \emph{Spaces.} $S$, $A$ finite; continuous or token-vocabulary-sized spaces are covered only to the extent that they are finite.
\end{itemize}

\paragraph{Reading the pieces together.}
What Lean does establish is the shared mechanism: a Markov product turns into a sum of per-step terms; recursions over the chain replace sums over exponentially many paths (forward and backward for the CRF, Viterbi, Bellman for the value function); and normalization of a local factor turns its score into a zero-mean term.
Which normalization one chooses decides whether the recursion computes a partition function (CRF), a probability that is already one (HMM, policy), or a maximum (Viterbi).
These are the objects that appear on the way from CRF to PPO.

\paragraph{Conclusion.}
We gave a Lean~4 account of three recursions on the path from linear-chain CRFs to policy-gradient reinforcement learning.
For the CRF, we formalized the Markov product, the log-score decomposition, the forward recursion with the negative log-likelihood, the Viterbi value recursion, and a forward--backward theorem for arbitrary weights with the backward recursion and the identity $\sum_{ys_t=a}P_m(ys)=z_t(a)\beta_t(a)$.
For reinforcement learning, we formalized the MDP trajectory measure, the Bellman equations, the policy gradient theorem in truncated, action-value and REINFORCE form, the zero-mean score, the unbiased advantage and GAE with the exact $V$.
We separated carefully what is a theorem from what is a conceptual comparison.
Natural next steps are the pairwise and normalized marginals and the identity $\grad\log Z=\E[\grad\,\mathrm{score}]$ for a general-potential CRF, a Viterbi argmax with back-pointers, a locally normalized autoregressive model as an MDP with deterministic transitions, and the importance-ratio identity behind the PPO surrogate.

\paragraph{Acknowledgements.}
Mathlib~\cite{mathlib2020,mathlib} provides the measure theory and analysis on which every theorem here rests.
Su Jianlin's articles on his blog Scientific Spaces were the guide to the CRF material~\cite{su2017crf,su2018search}.

\bibliographystyle{plain}
\bibliography{refs}

\medskip
\noindent\small Interactive statements of the modules cited in this paper are available at \href{http://www.lemma.cn/}{lemma.cn}; the project is on \href{https://github.com/math-proof/lemma}{GitHub}: \href{http://www.lemma.cn/lean/?module=Tensor.EqDot_GradExpect.of.Eq_Conditioned.Eq_Expect.IsFinite.IsFinite.generalized_advantage_estimate}{Lean~4}.

\end{document}
